\clearpage
\section{Renaming all lambda-free terms and padded binary clauses}
\label{sec:copying-lambda-free-renaming}
\termdef{copying-lambda-free-renaming}
The main theorem of this milestone proves bijective environment renaming for
every lambda-free term of the original syntax, at arbitrary finite variable and
atom scopes. Application, all unary operations, and all four binary operations
may be nested without restriction. Repeated variable or atom occurrences are
allowed: both two-use distributor forests and the original middle-block exchange
remain in the graph. A named distinct-list environment corollary uses exactly
the existing environment and translation definitions.

The construction also proves application and binary congruence for supplied
operand isomorphisms at arbitrary, possibly unequal allowances. Thus it is usable
with boxed operands once their own boundary correspondence is established. It
does not prove lambda-box renaming, an operational allocation/read/write rule,
or Proposition~12.2. The theorem at the end is stated on the original
\texttt{Translation.term}; no replacement translation is introduced.

\Needspace{12\baselineskip}
\subsection{Coordinates of the original operand shuffle}

\termdef{copying-pair-boundary-coordinates}
The left and right operand tags are represented by Booleans, with use numbers
zero and one. A component index pairs such a tag with an original environment
position. The grouped position is obtained by applying the existing pair-boundary
bijection and transporting through the existing interface-length equality.
Every grouped position has this form: invert the pair-boundary bijection after
transporting back along the length equality, and split its component tag.

\lean{Copying/PairBoundaryCoordinates}{8}{8}{PairBoundaryCoordinates.use}
\lean{Copying/PairBoundaryCoordinates}{10}{12}{PairBoundaryCoordinates.component}
\lean{Copying/PairBoundaryCoordinates}{15}{18}{PairBoundaryCoordinates.index}
\lean{Copying/PairBoundaryCoordinates}{21}{31}{PairBoundaryCoordinates.index_surjective}

The original shuffle is the inverse of the original four-block ordering
followed by the shuffled ordering. Evaluating it at an original component
therefore cancels the original ordering's inverse. Apply this cancellation in
the four cases (left/right operand, term/atom resource). Term resource $i$ in use
$u$ reaches grouped term position $i+ku$. Atom resource $a$ reaches grouped atom
position $a+hu$, after both term blocks; its absolute position is $2k+a+hu$.
For atoms, the separate concatenation identity checks that $k+a$ is exactly the
right-block injection, by equality of finite-index values. These calculations
retain the original order rather than replacing the shuffle by a counting law.

\lean{Copying/PairBoundaryCoordinates}{33}{39}{PairBoundaryCoordinates.shuffle_component}
\lean{Copying/PairBoundaryCoordinates}{42}{65}{PairBoundaryCoordinates.term_index}
\lean{Copying/PairBoundaryCoordinates}{67}{70}{PairBoundaryCoordinates.atom_join}
\lean{Copying/PairBoundaryCoordinates}{73}{97}{PairBoundaryCoordinates.atom_index}

\Needspace{12\baselineskip}
\subsection{Matching the multi-use environment forests}

\termdef{copying-multi-use-boundary-action}
Tensor the original automorphisms of $D_{k,m}$ and $D_{h,m}$. This constructs an
automorphism of the actual environment distributor for every use count $m$ and
allowance. Its outputs carry the resource actions $\rho$ and $\alpha$. Its input
actions preserve the use number: $i+ku$ maps to $\rho(i)+ku$, and $a+hu$ maps to
$\alpha(a)+hu$. The existing finite-grid index supplies the bounds even for
arbitrary $m$; the original distributor action and tensor block laws prove the
ordered-port equations. In particular, zero resources and zero uses need no
nonemptiness assumption.

\lean{Copying/MultiUseBoundaryRenaming}{8}{10}{MultiUseBoundaryRenaming.iso}
\lean{Copying/MultiUseBoundaryRenaming}{12}{17}{MultiUseBoundaryRenaming.output_term}
\lean{Copying/MultiUseBoundaryRenaming}{19}{24}{MultiUseBoundaryRenaming.output_atom}
\lean{Copying/MultiUseBoundaryRenaming}{26}{38}{MultiUseBoundaryRenaming.input_term}
\lean{Copying/MultiUseBoundaryRenaming}{40}{56}{MultiUseBoundaryRenaming.input_atom}

A typed environment permutation $\pi$ splits into its term and atom
bijections by the previous section. Apply the multi-use construction to that
pair. Splitting an arbitrary output position into the two resource blocks proves
that the output action is exactly $\pi$. For two uses, the coordinate laws above
identify the input action with applying $\pi$ independently inside each original
operand, then performing the original pair shuffle. This is the precise bridge
between the two operand graphs and their sharing forests.

\lean{Copying/MultiUseBoundaryRenaming}{58}{58}{MultiUseBoundaryRenaming.ofTyped}
\lean{Copying/MultiUseBoundaryRenaming}{61}{71}{MultiUseBoundaryRenaming.ofTyped_output}
\lean{Copying/MultiUseBoundaryRenaming}{75}{86}{MultiUseBoundaryRenaming.ofTyped_pair_input}

\Needspace{12\baselineskip}
\subsection{The application and binary constructor proofs}

A typed boundary isomorphism fixes a singleton input index: both its image
and the original index are zero. Assemble two operand isomorphisms by the existing
parallel theorem using the original pair boundary; transport the resulting
output interface along the same equality as the translation. This defines the
raw and final pair isomorphisms. Their two argument inputs are individually fixed,
as follows by the singleton law and the two component insertion equations.
If the operand output maps agree, both shuffled output blocks carry that common
map; prove this by the component output equation in each of the two cases.

\lean{Copying/PairBoundaryRenaming}{8}{14}{PairBoundaryRenaming.input_singleton}
\lean{Copying/PairBoundaryRenaming}{22}{26}{PairBoundaryRenaming.rawIso}
\lean{Copying/PairBoundaryRenaming}{28}{29}{PairBoundaryRenaming.iso}
\lean{Copying/PairBoundaryRenaming}{32}{57}{PairBoundaryRenaming.iso_input}
\lean{Copying/PairBoundaryRenaming}{60}{79}{PairBoundaryRenaming.iso_output}

A small head index selects either the original application edge or one
of the original binary-operation edges. In both cases the selected graph is the
existing primitive graph with one term input and two ordered term outputs. It
is an abbreviation for these original graphs, not an additional language
operation. The identity of this head matches the fixed two argument inputs of
the pair isomorphism, so serial gluing yields the prefix isomorphism.

\lean{Copying/BinaryBoundaryRenaming}{8}{8}{BinaryBoundaryRenaming.Head}
\lean{Copying/BinaryBoundaryRenaming}{10}{15}{BinaryBoundaryRenaming.graph}
\lean{Copying/BinaryBoundaryRenaming}{17}{21}{BinaryBoundaryRenaming.prefix_compatible}
\lean{Copying/BinaryBoundaryRenaming}{23}{23}{BinaryBoundaryRenaming.prefixIso}

Choose the actual two-use forest automorphism from the left operand's
typed output bijection. To prove compatibility, express every intermediate
position as a component and an environment position using surjectivity of the
pair index. The prefix output law, the pair output law, and the two-use input
law all give the same grouped position. Thus agreement between operand resource
actions proves the entire matching obligation. The clause abbreviation is
literally the original serial expression, and gluing this prefix and forest
gives its boundary isomorphism. The external output action equals the original
left operand action because the final forest realises that action exactly.

\lean{Copying/BinaryBoundaryRenaming}{25}{25}{BinaryBoundaryRenaming.forestIso}
\lean{Copying/BinaryBoundaryRenaming}{28}{38}{BinaryBoundaryRenaming.forest_compatible}
\lean{Copying/BinaryBoundaryRenaming}{40}{41}{BinaryBoundaryRenaming.clause}
\lean{Copying/BinaryBoundaryRenaming}{43}{45}{BinaryBoundaryRenaming.iso}
\lean{Copying/BinaryBoundaryRenaming}{47}{52}{BinaryBoundaryRenaming.iso_output}

Specialise the head to application or to each binary operation. This gives
public constructor isomorphisms for the existing translation functions, together
with their exact external output equations. The only compatibility hypothesis
at this constructor level is equality of the two operand output actions; below,
the syntax induction proves it from a common resource specification.

\lean{Copying/BinaryBoundaryRenaming}{61}{62}{BinaryBoundaryRenaming.application_boundaryIso}
\lean{Copying/BinaryBoundaryRenaming}{64}{65}{BinaryBoundaryRenaming.binary_boundaryIso}
\lean{Copying/BinaryBoundaryRenaming}{67}{69}{BinaryBoundaryRenaming.application_boundaryIso_output}
\lean{Copying/BinaryBoundaryRenaming}{71}{73}{BinaryBoundaryRenaming.binary_boundaryIso_output}

\Needspace{12\baselineskip}
\subsection{The resource-action induction invariant}

\termdef{copying-resource-boundary-action}
For resource bijections $\rho$ and $\alpha$, a resource-action certificate says
that an isomorphism maps variable output $i$ to $\rho(i)$ and atom output $k+a$
to $k+\alpha(a)$. Two such certificates imply equality of their complete output
index maps: split an arbitrary environment index at $k$, apply the relevant
certificate equation, and use injectivity of the output-port constructor.
This is why a mere existence theorem for each operand is too weak for induction.

\lean{Copying/ResourceBoundaryAction}{8}{11}{ResourceBoundaryAction.ResourceAction}
\lean{Copying/ResourceBoundaryAction}{18}{33}{ResourceBoundaryAction.agree}

If a constructor's external action is the operand's action, compose that
output equation with the operand certificate. This proves a general certificate
transfer lemma. Apply it to the existing unary output theorem and to the new
application and binary output theorems, using the agreement just proved for the
two operands. Consequently the same specified resource action persists through
every non-binding constructor.

\lean{Copying/ResourceBoundaryAction}{36}{46}{ResourceBoundaryAction.of_output}
\lean{Copying/ResourceBoundaryAction}{48}{50}{ResourceBoundaryAction.unary}
\lean{Copying/ResourceBoundaryAction}{52}{54}{ResourceBoundaryAction.application}
\lean{Copying/ResourceBoundaryAction}{56}{58}{ResourceBoundaryAction.binary}

The atomic isomorphisms must satisfy that stronger specification too.
Extract the selected-resource output index from its existing explicit port
equation. Tensor it with the unused atom forest to prove the variable clause's
two resource equations. For the instance prefix, insertion commutation at an
output says that its external output map is exactly the operand output map;
checking the allowance cases proves this for the original prefix. Combine it
with the selected-atom action and unused variable forest to prove the atom clause.
For constants, the primitive has no outputs, so the inner tensor's outputs all
come from the variable weakening forest and the outer tensor adds exactly the
atom block. Tensor output laws and the empty-output arithmetic prove both
constant equations. All existing interface casts are retained throughout.

\lean{Copying/AtomicBoundaryActions}{8}{14}{AtomicBoundaryActions.selected_outputIndex}
\lean{Copying/AtomicBoundaryActions}{16}{33}{AtomicBoundaryActions.variableTerm_action}
\lean{Copying/AtomicBoundaryActions}{39}{44}{AtomicBoundaryActions.iso_output}
\lean{Copying/AtomicBoundaryActions}{50}{70}{AtomicBoundaryActions.atom_action}
\lean{Copying/AtomicBoundaryActions}{72}{103}{AtomicBoundaryActions.constant_action}

\Needspace{12\baselineskip}
\subsection{Induction on every lambda-free term}

The lambda-free predicate is indexed by the original term. It includes
variables, atoms, constants, and is closed under every unary and binary operation
and application. It has no lambda constructor. Ordinary syntactic renaming
preserves membership by induction, including non-injective renamings; this
membership fact alone is not a graph substitution or non-injective graph-renaming
theorem. A second induction shows the translation allowance is zero: leaves
have depth zero, unary depth is unchanged, and the maximum of two zeros is zero.

\lean{Copying/LambdaFreeRenaming}{9}{15}{LambdaFreeRenaming.LambdaFree}
\lean{Copying/LambdaFreeRenaming}{19}{27}{LambdaFreeRenaming.rename}
\lean{Copying/LambdaFreeRenaming}{29}{30}{LambdaFreeRenaming.depth_zero}

For the induction only, expose this zero allowance by equality transport
of the already defined translated graph. This auxiliary graph is the original
translation with its index rewritten, not a simplification of its carriers or
edges. Equality transport commutes with unary translation. For application and
binary translation, eliminate the two depth-zero equalities: their maximum is
zero and both original pads compute to the zero pad. This proves the two general
zero-allowance identities and their specialisations to actual translated terms.

\lean{Copying/LambdaFreeRenaming}{34}{35}{LambdaFreeRenaming.graph}
\lean{Copying/LambdaFreeRenaming}{37}{39}{LambdaFreeRenaming.graph_unary}
\lean{Copying/LambdaFreeRenaming}{42}{48}{LambdaFreeRenaming.application_zero}
\lean{Copying/LambdaFreeRenaming}{50}{56}{LambdaFreeRenaming.binary_zero}
\lean{Copying/LambdaFreeRenaming}{58}{60}{LambdaFreeRenaming.graph_application}
\lean{Copying/LambdaFreeRenaming}{62}{64}{LambdaFreeRenaming.graph_binary}

The strong induction theorem constructs an isomorphism together with its
resource-action certificate. Leaves use the newly certified atomic isomorphisms.
A unary node uses unary constructor congruence and certificate preservation.
At application or binary nodes, obtain both certified operand isomorphisms,
derive agreement from the certificates, apply constructor congruence, and carry
the certificate forward. Thus the agreement premise is discharged inside the
proof, at every node and for arbitrarily repeated resources.

Finally, eliminating the two zero-allowance equalities converts the induction's
isomorphism back to the original translated graphs. The theorem
\texttt{Translation.lambdaFree\_rename} has only the original term, its
lambda-free certificate, and the two finite resource bijections as parameters.
It has no operational, stability, robustness, or operand-agreement assumption.
The only target equality transport is the pre-existing syntactic depth-renaming
identity, just as in the earlier atomic endpoint.

\lean{Copying/LambdaFreeRenaming}{68}{91}{LambdaFreeRenaming.rename_with_action}
\lean{Copying/LambdaFreeRenaming}{94}{98}{LambdaFreeRenaming.uncast}
\lean{Copying/LambdaFreeRenaming}{104}{109}{LambdaFreeRenaming.lambdaFree_rename}

\Needspace{12\baselineskip}
\subsection{Unequal allowances and supplied boxed operands}

\termdef{copying-boundary-padding-renaming}
Recursive label equivalence survives one step of the original padding. Primitive
labels retain equality. For boxes, their existing body isomorphism is padded by
the previously proved recursive-isomorphism theorem, preserving binder metadata.
This yields a padded boundary isomorphism with the exact same outer port and edge
bijections. Its external output action is unchanged. Iteration gives any finite
padding, with the output equation proved by induction on the number of steps.
These are structural representation results, not execution-padding invariance.

\lean{Copying/BoundaryPaddingRenaming}{8}{23}{BoundaryPaddingRenaming.label_equiv}
\lean{Copying/BoundaryPaddingRenaming}{31}{46}{BoundaryPaddingRenaming.raise}
\lean{Copying/BoundaryPaddingRenaming}{48}{50}{BoundaryPaddingRenaming.raise_output}
\lean{Copying/BoundaryPaddingRenaming}{52}{54}{BoundaryPaddingRenaming.pad}
\lean{Copying/BoundaryPaddingRenaming}{56}{63}{BoundaryPaddingRenaming.pad_output}

Equality transport of the allowance also preserves the isomorphism and
its output action; eliminate the equality to prove both facts. Compose iterated
padding and equality transport to obtain the original \texttt{padTo} operation,
whose output law is obtained by composing the two preceding laws. Substituting
that equation into the term and atom certificate fields proves resource-action
preservation under padding.

\lean{Copying/BoundaryPaddingRenaming}{65}{67}{BoundaryPaddingRenaming.castDepth}
\lean{Copying/BoundaryPaddingRenaming}{69}{71}{BoundaryPaddingRenaming.castDepth_output}
\lean{Copying/BoundaryPaddingRenaming}{73}{75}{BoundaryIso.padTo}
\lean{Copying/BoundaryPaddingRenaming}{77}{82}{BoundaryPaddingRenaming.padTo_output}
\lean{Copying/BoundaryPaddingRenaming}{88}{99}{ResourceAction.padTo}

For operands at depths $d$ and $e$, pad their supplied certified
isomorphisms independently to $\max(d,e)$. Their resource certificates still
agree, so the application and binary constructor theorems apply at exactly the
common allowance chosen by Figure~16. No equal-depth premise is added, and no
box is removed. These two constructors are ready for operands containing boxes,
provided their own resource-renaming isomorphisms are established separately.

\lean{Copying/BoundaryPaddingRenaming}{111}{115}{BoundaryPaddingRenaming.padded_application_boundaryIso}
\lean{Copying/BoundaryPaddingRenaming}{117}{121}{BoundaryPaddingRenaming.padded_binary_boundaryIso}

\Needspace{12\baselineskip}
\subsection{Concrete instances, rejected mismatch, and remaining scope}

The callback invocation $f()$ is an application of a variable to the
original unit constant, so its lambda-free certificate immediately gives the
renaming theorem. Self-application of a variable gives a regression with the
same resource in both operands. Assignment of an allocated reference combines
the original atom, reference-creation, constant and assignment clauses; the
same endpoint proves its renaming. A further family combines dereferencing a
variable with an atom and is parameterised by every binary operation, covering
addition, subtraction, assignment and atom equality without separate assumptions.
These are translation theorems, not termination or store-behaviour assertions.

\lean{Copying/LambdaFreeRenamingExamples}{8}{12}{LambdaFreeRenamingExamples.callbackInvocation_rename}
\lean{Copying/LambdaFreeRenamingExamples}{15}{19}{LambdaFreeRenamingExamples.repeatedVariable_rename}
\lean{Copying/LambdaFreeRenamingExamples}{22}{26}{LambdaFreeRenamingExamples.assignmentAllocation_rename}
\lean{Copying/LambdaFreeRenamingExamples}{29}{33}{LambdaFreeRenamingExamples.everyBinary_rename}

A concrete two-variable, two-atom term applies a dereferenced variable
to an atom-equality expression. Its membership proof follows the syntax, and
swapping both resource sorts instantiates the general theorem nontrivially.
For the negative regression, compare two original constant-graph automorphisms:
one swaps two variable resources and one fixes them. Their certified images of
variable output zero differ. Equality of their whole output actions would imply
$1=0$. This rejects the tempting but invalid claim that well-typed operand
isomorphisms automatically match.

\lean{Copying/LambdaFreeRenamingExamples}{36}{37}{LambdaFreeRenamingExamples.twoSortExample}
\lean{Copying/LambdaFreeRenamingExamples}{39}{40}{LambdaFreeRenamingExamples.twoSortExample_free}
\lean{Copying/LambdaFreeRenamingExamples}{42}{45}{LambdaFreeRenamingExamples.twoSortExample_rename}
\lean{Copying/LambdaFreeRenamingExamples}{49}{58}{LambdaFreeRenamingExamples.mismatched_operand_actions}

Identity isomorphisms satisfy the identity resource certificate by
reflexivity. Further regressions use the proved nontrivial atom and variable
permutations, first padded to allowances one and three, then combined by the
common-allowance application and binary theorems. A genuine translated lambda
operand at depth one can also be combined with a depth-zero constant, using the
lambda operand's identity isomorphism. This last instance tests boxed composition,
not arbitrary lambda renaming. Finally the named-environment corollary specialises
the lambda-free theorem to the lengths of an existing distinct-list environment;
the environment translation wrapper is definitionally the original translation.

\lean{Copying/BinaryRenamingRegression}{7}{9}{BinaryRenamingRegression.refl}
\lean{Copying/BinaryRenamingRegression}{16}{19}{BinaryRenamingRegression.unequalAllowanceRenaming}
\lean{Copying/BinaryRenamingRegression}{22}{25}{BinaryRenamingRegression.unequalAllowanceBinaryRenaming}
\lean{Copying/BinaryRenamingRegression}{29}{32}{BinaryRenamingRegression.boxedOperandApplication}
\lean{Copying/BinaryRenamingRegression}{35}{42}{BinaryRenamingRegression.lambdaFree_namedEnvironment}

\paragraph{Completion boundary and estimate.}
Application and binary operations are removed from the remaining bijective
translation-renaming cases. The missing constructor is lambda, including the
correspondence between a box's external free-resource interface and its body's
ordered boundary. The existing recursive label relation fixes body boundaries;
it cannot silently serve as a proof of arbitrary boundary permutation inside a
box. Establish that correspondence while preserving the original terms and
Figure~16 translation. Non-bijective graph renaming, graph substitution, the
concrete copying UAM, its execution representation laws, all eleven robustness
template families, semantic weakening, and Figure~22's equivalence chain remain.

As of 25 October 2026, estimated completion of the requested Proposition~12.2
formalisation is approximately \textbf{50\% by mathematical effort}, with a broad
plausible range of \textbf{35--60\%}. This estimate supersedes the earlier
unary-fragment estimate. It reflects a complete non-binding syntax induction and
unequal-allowance constructor proofs; it does not count these as completed machine
or robustness-template proofs. Proposition~12.2 itself is not yet proved.
